%% kapa: exam=1 round=34 year=2023 period=1 form=A subject=civil_law subject_ko=민법 seq=5 number=5
%% source: Q-Net 감정평가사 기출문제 (cst00309) artlSeq=5211789 file=raw/1008/5211789/2023년도 제34회 감정평가사 1차시험 1교시.pdf
%% conversion: kordoc 4.24.51 -> tools/gen1.py (auto); review=proofread-sonnet; flags=none
\begin{question}{5}
비법인사단에 관한 설명으로 옳지 않은 것은? (다툼이 있으면 판례에 따름)
\begin{choices}
  \item 비법인사단의 대표자는 자신의 업무를 타인에게 포괄적으로 위임할 수 있다.
  \item 정관이나 규약에 달리 정함이 없는 한, 사원총회의 결의를 거치지 않은 총유물의 관리행위는 무효이다.
  \item 고유한 의미의 종중은 종중원의 신분이나 지위를 박탈할 수 없고, 종중원도 종중을 탈퇴할 수 없다.
  \item 고유한 의미의 종중은 자연발생적 종족단체이므로 특별한 조직행위나 성문의 규약을 필요로 하지 않는다.
  \item 비법인사단의 사원이 집합체로서 물건을 소유할 때에는 총유로 한다.
\end{choices}
\end{question}
\answer{1}
